\subsection{Locally trivial fiber spaces}

Let B be a topological space. If F is a topological space, we call the trivial B-fibered space with fiber type F the B-space \((B \times F, pr_{1})\) .

Let E be a B-space. If there exists a topological space F and a B-isomorphism \(u: E \rightarrow B \times F\), we say that the B-space E is a trivializable B-fibered space and that u is a trivialization of it with fiber type F.

DEFINITION 1. --- Let B be a topological space. Let E be a B-space and let p be its projection. One says that E is a locally trivial B-fibered space if every point of B has a neighborhood V such that \((\overline{p}^{1}(V), p_{V})\) is a trivializable V-fibered space.

Instead of “locally trivial fibered B-space,” one also says “locally trivial fibered space with base B.” If E is a locally trivial fibered B-space, one sometimes says that \((E,B,p)\) is a locally trivial fibration, or, by abuse of language, that p is a locally trivial fibration. One also says that a B-space \((E,p)\) is trivializable over a subset A of B if the A-space \(E_{A}\) induced by \((E,p)\) over A is a trivializable fiber space.

Let E be a locally trivial B-fibered space; denote its projection by p. The set of points a in B for which the fiber \(\overline{p}^{1}(a)\) is empty (resp. nonempty) is open. The image of p is therefore an open and closed subset of B.

Let F be a topological space. If all the fibers of E are homeomorphic to F, then E is said to be a locally trivial B-fibered space with fiber type F.

Remarks. --- 1) Let (E, p) be a locally trivial fiber bundle over B, and let A be a subset of B. The A-space \(\mathrm{E}_{\mathrm{A}} = (\stackrel{-1}{p}(\mathrm{A}), p_{\mathrm{A}})\) deduced from E by passing to subspaces is a locally trivial fiber space. If E is trivializable, then so is \(E_{A}\) .

Indeed, for every point \(a\) of A, there exists an open neighborhood U of \(a\) in B, a topological space F and a U-isomorphism \(g\colon \overline{p}^{1}(\mathrm{U})\to \mathrm{U}\times \mathrm{F}\) . The map \(g\) induces an \((\mathrm{A}\cap \mathrm{U})\) -isomorphism of \(\overline{p}_{\mathrm{A}}^{-1}(\mathrm{A}\cap \mathrm{U})\)

onto \((\mathrm{A}\cap\mathrm{U})\times\mathrm{F}\) , which proves that \(E_{A}\) is a locally trivial fibered A-space.

\begin{enumerate}
\def\labelenumi{\arabic{enumi})}
\setcounter{enumi}{1}
\tightlist
\item
  Let \((\mathrm{E}, p)\) be a B-space and let \((\mathrm{E}_{i})_{i \in \mathrm{I}}\) be a partition of E consisting of open sets.
\end{enumerate}

If each of the B-spaces \((\mathrm{E}_{i}, p \mid \mathrm{E}_{i})\) is a trivializable fiber space, E is a trivializable B-fibered space. Indeed, for each \(i \in I\), let \(F_{i}\) be a topological space such that the B-spaces \((\mathrm{E}_{i}, p \mid \mathrm{E}_{i})\) and \((\mathrm{B} \times \mathrm{F}_{i}, \mathrm{pr}_{1})\) are isomorphic. Then the B-space \((\mathrm{E}, p)\) is isomorphic to \((\mathrm{B} \times \mathrm{F}, \mathrm{pr}_{1})\), where F is the topological sum space of the family \((\mathrm{F}_{i})_{i \in I}\) .

Suppose the set I is finite. If each of the B-spaces \((\mathrm{E}_{i}, p \mid \mathrm{E}_{i})\) is a locally trivial B-fibered space, the same holds for the B-space E. Indeed, since the set I is finite, each point of B has a neighborhood U over which the B-fibered spaces \(E_{i}\) are all trivializable. By the preceding, the U-space \((\bar{p}^{1}(\mathrm{U}), p_{\mathrm{U}})\) is then a trivializable fiber bundle.

